\chapter{Metode Penelitian}

% ================================================================
% PANDUAN PENULISAN BAB III — METODE PENELITIAN
% Bab ini menjelaskan cara penelitian dilakukan secara rinci sehingga
% dapat direplikasi oleh peneliti lain. Tuliskan secara sistematis
% mengikuti urutan bagian berikut.
% ================================================================

\section{Desain Penelitian}

% Sebutkan desain penelitian Anda, misalnya deskriptif (cross-sectional
% atau longitudinal), case-control, kohort, atau eksperimental. Jelaskan
% singkat alasan pemilihan desain tersebut.

\section{Tempat dan Waktu Penelitian}

% Cantumkan lokasi pengambilan data (misal unit/instalasi rumah sakit,
% klinik, atau laboratorium) serta periode/rentang waktu pengumpulan data.

\section{Populasi dan Sampel Penelitian}

% - Bedakan populasi target dan populasi terjangkau.
% - Jelaskan teknik pengambilan sampel (total sampling, consecutive
%   sampling, purposive sampling, dll).
% - Cantumkan perhitungan besar sampel bila diperlukan.

\subsection{Kriteria Inklusi}

% Daftar karakteristik subjek yang memenuhi syarat untuk diteliti.
% \begin{itemize}
%     \item Kriteria inklusi pertama.
%     \item Kriteria inklusi kedua.
% \end{itemize}

\subsection{Kriteria Eksklusi}

% Daftar karakteristik yang menyingkirkan subjek dari penelitian.
% \begin{itemize}
%     \item Kriteria eksklusi pertama.
%     \item Kriteria eksklusi kedua.
% \end{itemize}

\section{Definisi Operasional}

% Sajikan definisi operasional variabel dalam bentuk tabel: nama
% variabel, definisi, alat ukur, cara ukur, dan hasil ukur/skala.
% Contoh format tabel:
%
% \begin{table}[ht]
% \centering
% \begin{tabular}{p{3.2cm}p{3.6cm}p{2.4cm}p{3cm}p{2.4cm}}
% \toprule
% Variabel & Definisi & Alat Ukur & Cara Ukur & Hasil/Skala \\
% \midrule
% Variabel 1 & ... & ... & ... & ... \\
% Variabel 2 & ... & ... & ... & ... \\
% \bottomrule
% \end{tabular}
% \caption{Definisi operasional variabel}
% \label{tab:definisi}
% \end{table}

\section{Instrumen Penelitian}

% Jelaskan instrumen/alat yang digunakan (kuesioner, formulir pengumpulan
% data, alat pemeriksaan, perangkat lunak, dll) beserta validitas dan
% reliabilitasnya bila relevan.

\section{Prosedur Pengumpulan Data}

% Uraikan langkah pengumpulan data secara berurutan: persiapan,
% seleksi subjek, pemeriksaan/pengisian instrumen, hingga rekapitulasi
% dan verifikasi data.

\section{Analisis Data}

% Sebutkan jenis analisis yang digunakan (deskriptif dan/atau
% inferensial) serta perangkat lunak statistik. Untuk analisis
% inferensial, jelaskan uji statistik yang dipakai sesuai karakteristik
% data (skala, normalitas, distribusi).

\section{Etika Penelitian}

% Jelaskan persetujuan etik (nomor kode etik bila sudah didapat),
% prosedur informed consent, serta jaminan kerahasiaan data subjek.
